\begin{itemize}
\item \textbf{Simulated correctness and collection costs.} Verdicts come from benchmark gold answers. Flipped, missing and complaint-skewed verdicts are simulations, not measurements of human behaviour. Real feedback studies~\citep{kreutzer2018ebay,gao2023continually} motivate testing systematic rater errors and selection bias. Noise is tested only under repeated feedback, not the one-rating setting. Additional calibration and switch verdicts mean the pipeline is not a one-verdict total annotation budget. We do not measure live exploration cost, user outcomes or annotation time.
\item \textbf{Seeds and incomplete comparisons.} Table~\ref{tab:seeds} lists the available independent training runs. The main Vega feedback and labelled arms have three seeds per benchmark. Laya encoder arms have three on the public menu and JevBench, two on DecisionBench and one on typed-decisions and emotion; AG News and a matched one-pass labelled encoder arm were not run. The realistic-feedback arms have three seeds on DecisionBench and one on typed-decisions and emotion. No-KL arms have three on DecisionBench and Vega 0.8B JevBench, and one elsewhere. Head-only Laya one-pass arms have three on typed-decisions and one on the public menu, JevBench and DecisionBench. Historical head-only eight-pass results are identified separately; their labelled head lacked the KL term. Single-seed intervals cannot estimate variation across training runs. The failed half-rated emotion run at 4B is retained and disclosed in Section~\ref{sec:realistic}.
\item \textbf{Model and modality scope.} The study covers two released model families from one developer. Every evaluated task uses the text input path, including structured states rendered as text. Vega's image capability is not evaluated. A frozen vision encoder does not guarantee preserved image decisions when a shared downstream adapter changes; neither multimodal improvement nor retention is established. JevBench and DecisionBench are small, and the released models are already near 0.93 on AG News.
\item \textbf{Benchmark artifacts and overlap.} Five tasks have answer-position, option-structure or class-prior shortcuts. Vega's shipped typed-decisions adapter was fitted on the feedback pool; the engines' original training corpora are undisclosed. Removing the identified tasks and disabling shipped adapters are sensitivity analyses, not proof that every remaining task is free of contamination.
\item \textbf{Historical design choices.} The KL coefficient was selected in a pilot involving DecisionBench test rows (Appendix~\ref{app:pilot}), so that benchmark is not an untouched test of the full research process. Rank, learning rates and pass count were fixed before this grid. Labelled arms use calibration labels for early stopping; the main feedback arms do not receive that model-selection signal. The grid was reduced while running to fit the GPU budget (Appendix~\ref{app:repro}).
\item \textbf{Calibration and statistical precision.} One temperature per benchmark may be too coarse for a heterogeneous task menu. ECE is sensitive to binning and sample size, so proper scores are also reported. Non-significant differences are inconclusive rather than evidence of equality; the exploratory non-inferiority marks use unadjusted intervals. Three seeds give limited information about training variability, and single-seed intervals are conditional on the observed run.
\item \textbf{Reproduction and causal attribution.} Historical same-seed reruns differ across machines by up to about 3 points. The complete current grid has not been independently reproduced. Comparing logged and fresh on-policy feedback changes both information and the learning procedure; it does not identify a pure causal effect of verdict count.
\end{itemize}
